\documentclass[11pt]{article}

\usepackage[utf8]{inputenc}
\usepackage[T1]{fontenc}
\usepackage{lmodern}
\usepackage{geometry}
\usepackage{amsmath,amssymb,amsfonts,amsthm,mathtools}
\usepackage{bm}
\usepackage{enumitem}
\usepackage{microtype}
\usepackage{hyperref}
\usepackage{titlesec}
\usepackage{booktabs}
\usepackage{array}
\usepackage{setspace}
\usepackage{mathrsfs}
\usepackage{physics}

\geometry{margin=1in}
\hypersetup{colorlinks=true, linkcolor=black, urlcolor=blue, citecolor=black}
\setlist[enumerate]{topsep=0.4em,itemsep=0.35em}
\setlist[itemize]{topsep=0.3em,itemsep=0.25em}
\titleformat{\section}{\large\bfseries}{\thesection.}{0.5em}{}
\titleformat{\subsection}{\normalsize\bfseries}{\thesubsection.}{0.5em}{}
\setstretch{1.06}

\newcommand{\Aether}{\AE{}ther}
\newcommand{\Aflow}{\AE{}ther-flow}
\newcommand{\TheoryName}{\Aflow{} Interpretation of Relativity}
\newcommand{\TheoryProgram}{\Aether{} / \Aflow{} framework}
\newcommand{\Sub}{\mathcal{A}}
\newcommand{\BridgeMetric}{\mathfrak{g}}

\newtheorem{definition}{Definition}
\newtheorem{proposition}{Proposition}
\newtheorem{remark}{Remark}
\newtheorem{corollary}{Corollary}

\title{\textbf{The \TheoryName{}}\\[0.4em]
\large Next Candidate Family: Constraint-Assisted Longitudinal \texorpdfstring{\(U\)}{U}-Sector Completion}
\author{}
\date{}

\begin{document}

\maketitle

\begin{abstract}
The derivative-mixing \(S/Q\) branch has now been tested through its flat-background quadratic consistency pass and found to admit exact closure only on a tuned degenerate branch. The next derivational move must therefore change what both the minimal and derivative-mixing branches left untouched: the longitudinal \(U\)-sector constraint structure.

This note defines the next candidate family accordingly. It introduces an auxiliary scalar \(C\) that couples the divergence of the ordered-flow field to a sourced scalar combination of the order and \(Q\) sectors,
\[
\Delta\mathcal L_C
=
C\Bigl(\nabla_AU^A-\alpha\bigl(U^B\partial_BS-\sigma_0\bigr)-\beta_IQ^I\Bigr)
-\frac{M_C^2}{2}C^2.
\]
The bridge metric, the exact-closure benchmark, and universal matter coupling are left unchanged. On the flat admissible branch this generates the quadratic longitudinal block
\[
\Delta\mathcal L_C^{(2)}
=
C\Bigl(\nabla^2\chi+\frac{1}{2}\partial_0H-\alpha(\partial_0\sigma-\sigma_0\phi)-\beta_Iq^I\Bigr)
-\frac{M_C^2}{2}C^2,
\]
so the longitudinal \(U\)-sector is no longer eliminated by a vanishing bracket. Instead it becomes a sourced auxiliary channel feeding back into the scalar exact-closure problem.

The present note does not claim that this branch already succeeds. Its claim is narrower. It identifies the next candidate family worth testing after the failure of the derivative-mixing branch, writes its action explicitly, and shows why the previous no-go decompositions no longer apply without modification.
\end{abstract}

\section{Introduction}

The current manuscript sequence of \TheoryName{} already contains the exact-closure theory statement, its effective dynamics, its consistency analysis, its relativistic recovery statement, and the explicit substrate-bridge line. On the derivational side, the project has now passed through three increasingly sharp filters. First, the minimal candidate action was written explicitly. Second, the coefficient-level linearized consistency manuscript compressed its flat-background exact-closure burden into the scalar operator \(\mathcal B(k^2)\). Third, the minimal stable-sign no-go note and the derivative-mixing branch consistency test showed that neither the minimal branch nor the first nonminimal \(S/Q\) derivative-mixing branch supplies the sought nondegenerate flat-background exact-closure solution.

That outcome narrows the next move considerably. Both of the previous branches left the longitudinal \(U\)-sector structurally trivial in the same decisive sense: after variation, the \(\chi\) equation still reduced to a vanishing bracket and contributed no independent residual bridge self-energy. The next candidate family therefore has to alter that fact rather than merely reweighting the already-tested scalar operator downstream.

\paragraph{Framework and claim boundary.}
\TheoryProgram{} denotes the broader \Aether{} / \Aflow{} framework built on the claims that the \Aether{} is the underlying four-dimensional substrate of reality, the \Aflow{} is its intrinsic ordered motion, observed three-dimensional space is the local experiential slice of that deeper substrate, S-time is the experienced order of change arising from matter, light, and the \Aflow{}, and observed expansion is the three-dimensional appearance of deeper four-dimensional ordered motion. Within that framework, \TheoryName{} denotes the exact-closure theory in which the effective gravitational dynamics are adopted to be exactly Einsteinian with universal matter coupling. In the active sequence, \emph{adoption} means use of that established relativistic dynamics without claiming substrate derivation, while \emph{derivation} is reserved for a first-principles recovery from explicit substrate variables.

The present note stays inside that boundary. It does not claim a completed Einsteinian derivation. It nominates and formulates the next candidate family worth testing.

\section{Why the Next Branch Must Activate the Longitudinal \texorpdfstring{\(U\)}{U}-Sector}

The minimal and derivative-mixing branches share one decisive structural limitation. In both cases the longitudinal \(U\)-sector scalar
\[
\Theta \equiv \nabla_AU^A
\quad\leadsto\quad
\nabla^2\chi+\frac{1}{2}\partial_0H
\]
entered only through a quadratic term that vanished on its own decaying on-shell branch. As a result:
\begin{enumerate}
    \item the longitudinal \(U\)-sector did not remain as an independent sourced auxiliary channel;
    \item the exact-closure burden collapsed back to the \((\sigma,q^I)\) scalar block alone;
    \item every subsequent no-go result was decided without any genuine new longitudinal \(U\)-sector feedback.
\end{enumerate}

\begin{proposition}[Selection principle after the derivative-mixing result]
The next candidate family worth testing must modify the longitudinal \(U\)-sector constraint structure directly rather than only reweighting the already-tested scalar operator of the \(S/Q\) sector.
\end{proposition}

\begin{proof}
The derivative-mixing consistency manuscript showed that changing only the \(S/Q\) block still leaves the \(\chi\) equation structurally trivial and therefore fails to produce a nondegenerate exact-closure branch. A new family must therefore target what remained unchanged, namely the longitudinal \(U\)-sector elimination itself.
\end{proof}

\section{Constraint-Assisted Longitudinal \texorpdfstring{\(U\)}{U}-Sector Family}

\begin{definition}[Constraint-assisted \(U\)-sector completion]
The constraint-assisted longitudinal \(U\)-sector family is obtained by adjoining to the auxiliary sector a new scalar \(C\) with interaction
\begin{equation}
\Delta\mathcal L_C
=
C\Bigl(\nabla_AU^A-\alpha\bigl(U^B\partial_BS-\sigma_0\bigr)-\beta_IQ^I\Bigr)
-\frac{M_C^2}{2}C^2,
\label{eq:LC}
\end{equation}
where \(\alpha\), \(\beta_I\), and \(M_C^2\) are branch parameters and \(C\) carries no independent kinetic term at this stage.
\end{definition}

The full auxiliary sector of the new family is therefore
\begin{equation}
\mathcal L_{\mathrm{aux}}^{(C)}
=
\mathcal L_{\mathrm{aux}}
+\Delta\mathcal L_C.
\label{eq:LauxC}
\end{equation}

\begin{remark}
This is a genuine \(U\)-sector completion rather than another scalar-only deformation. The new field \(C\) is introduced only to mediate a sourced constraint on \(\nabla_AU^A\), so the modification is concentrated exactly where the previous branches remained inert.
\end{remark}

\subsection{Benchmark Structures Left Unchanged}

Three structural pieces are intentionally not changed.
\begin{enumerate}
    \item The bridge metric remains
    \begin{equation}
    \BridgeMetric_{AB}=G_{AB}-2U_AU_B.
    \label{eq:bridgemetric}
    \end{equation}
    \item The observer-accessible matter route remains \(S_{\mathrm{matter}}[\Xi^*\BridgeMetric,\psi]\), so the family still has a single universal matter metric.
    \item The exact-closure benchmark of \TheoryName{} remains the same Einsteinian observer-accessible limit already fixed elsewhere in the active sequence.
\end{enumerate}

Thus the new family changes the derivational architecture without changing the public theory claim.

\subsection{Reversion Limit}

The intended exact-closure reversion route is any limit in which the new constraint mediator decouples, for example
\begin{equation}
\alpha\to 0,
\qquad
\beta_I\to 0,
\label{eq:reversion1}
\end{equation}
or, equivalently at fixed source coefficients,
\begin{equation}
M_C^2\to\infty.
\label{eq:reversion2}
\end{equation}
In those limits the family contracts back to the previously studied candidate action line.

\section{Flat-Background Quadratic Structure}

Use the same admissible homogeneous background as in the previous flat-background tests:
\begin{equation}
\bar G_{AB}=\delta_{AB},
\qquad
\bar U^A=\delta^A{}_0,
\qquad
\bar S=\sigma_0X^0,
\qquad
\bar Q^I=0,
\label{eq:background}
\end{equation}
with perturbations
\begin{equation}
G_{AB}=\delta_{AB}+H_{AB},
\qquad
U^A=\bar U^A+V^A,
\qquad
S=\sigma_0X^0+\sigma,
\qquad
Q^I=q^I.
\label{eq:perturbations}
\end{equation}
As before,
\begin{equation}
V_i=V_i^T+\partial_i\chi,
\qquad
\partial_iV_T^i=0.
\label{eq:Vdecomp}
\end{equation}

At quadratic order the new interaction becomes
\begin{equation}
\Delta\mathcal L_C^{(2)}
=
C\Bigl(\nabla^2\chi+\frac{1}{2}\partial_0H-\alpha(\partial_0\sigma-\sigma_0\phi)-\beta_Iq^I\Bigr)
-\frac{M_C^2}{2}C^2,
\label{eq:LC2}
\end{equation}
where \(h_{00}=-2\phi\) and \(H=\delta^{ij}h_{ij}\).

Define the standard longitudinal scalar and order-sector combination
\begin{equation}
\Theta \equiv \nabla^2\chi+\frac{1}{2}\partial_0H,
\qquad
X \equiv \partial_0\sigma-\sigma_0\phi.
\label{eq:ThetaX}
\end{equation}
Then the full longitudinal block becomes
\begin{equation}
\mathcal L_{\Theta C}^{(2)}
=
-\frac{\kappa_\theta}{2}\Theta^2
+C\bigl(\Theta-\alpha X-\beta_Iq^I\bigr)
-\frac{M_C^2}{2}C^2.
\label{eq:LThetaC}
\end{equation}

\begin{proposition}[Why the previous no-go decompositions do not transfer directly]
On the constraint-assisted \(U\)-sector family, the longitudinal \(U\)-sector equations no longer imply \(\Theta=0\) on the decaying flat branch. Instead they imply a sourced auxiliary relation between \(\Theta\), \(X\), \(q^I\), and \(C\).
\end{proposition}

\begin{proof}
Varying \eqref{eq:LThetaC} with respect to \(C\) and \(\chi\) gives
\begin{equation}
\Theta-\alpha X-\beta_Iq^I-M_C^2C=0,
\label{eq:Ceq}
\end{equation}
and
\begin{equation}
\nabla^2\bigl(-\kappa_\theta\Theta+C\bigr)=0.
\label{eq:chieqC}
\end{equation}
On the decaying branch, \eqref{eq:chieqC} implies
\begin{equation}
C=\kappa_\theta\Theta.
\label{eq:CTheta}
\end{equation}
Substituting \eqref{eq:CTheta} into \eqref{eq:Ceq} yields
\begin{equation}
\bigl(1-\kappa_\theta M_C^2\bigr)\Theta
=
\alpha X+\beta_Iq^I.
\label{eq:sourcedTheta}
\end{equation}
Therefore \(\Theta\) is sourced by the scalar order and \(Q\) sectors rather than forced to vanish identically.
\end{proof}

\begin{corollary}
The flat-background exact-closure problem of the constraint-assisted \(U\)-sector family is not determined solely by the previously studied scalar operator \(\mathcal B(k^2)\) or its derivative-mixing generalization. The longitudinal \(U\)-sector must now be eliminated together with the \((\sigma,q^I)\) block as one coupled auxiliary system.
\end{corollary}

\section{Why This Is the Correct Next Family to Test}

The constraint-assisted family is the right next candidate family for four reasons.
\begin{enumerate}
    \item It targets exactly the sector that remained inert in the minimal and derivative-mixing branches.
    \item It preserves the bridge metric, the single-metric matter route, and the exact-closure benchmark.
    \item It introduces only one new auxiliary scalar and one sourced constraint relation, so the conceptual drift is still small.
    \item It is immediately testable by the same Phase 3 method already used on the earlier branches: derive the full quadratic auxiliary system, integrate out the coupled \((\chi,C,q^I,\sigma)\) block, and determine the resulting bridge self-energy.
\end{enumerate}

\begin{remark}
The family is not yet being promoted as viable. It is being promoted as the next place where a genuinely new cancellation mechanism could live.
\end{remark}

\section{Immediate Research Tasks for This Family}

Once the family is fixed by \eqref{eq:LC}, the next calculations are specific.
\begin{enumerate}
    \item derive the full quadratic auxiliary Lagrangian including the new \(C\)-sector
    \item eliminate the coupled \((\chi,C)\) system exactly
    \item determine the generalized scalar operator that replaces the previous \(\mathcal B(k^2)\)-only description
    \item check whether any nondegenerate flat-background exact-closure branch survives
    \item impose the corresponding health conditions on the sourced longitudinal block
\end{enumerate}

Two immediate parameter cautions also have to be tracked.
\begin{enumerate}
    \item The sourced longitudinal relation \eqref{eq:sourcedTheta} is singular at
    \begin{equation}
    1-\kappa_\theta M_C^2=0.
    \label{eq:singular}
    \end{equation}
    That locus must be treated as a separate degenerate branch, not as part of the generic family.
    \item The effective sign of the sourced longitudinal block must be checked explicitly after eliminating \(C\) and \(\chi\); it cannot be read off safely from \eqref{eq:LC} alone.
\end{enumerate}

\section{Conclusion}

The next derivational move after the derivative-mixing result had to be more than another scalar-sector retuning. It had to activate the longitudinal \(U\)-sector that both previous branches left structurally trivial. This note defines exactly that next candidate family.

The constraint-assisted longitudinal \(U\)-sector completion introduces one auxiliary scalar \(C\) that sources the divergence of the ordered-flow field by the order and \(Q\) sectors. As a result, the \(\chi\) equation no longer collapses to a vanishing bracket. The longitudinal \(U\)-sector becomes a real sourced auxiliary channel, so the previous no-go decompositions no longer transfer unchanged.

That is enough to make this the correct next family to test. It is not yet enough to claim success. The next manuscript has to perform the full quadratic elimination and determine whether the sourced longitudinal channel yields a genuinely nondegenerate flat-background exact-closure branch or only another disguised degeneration. Until that calculation is done, exact closure remains the benchmark and reversion theory of the project.

\end{document}
